\section{Variational setting, dictionaries, and algorithms}
\label{sec:setting}

Let \(X\) be a real reflexive Banach space, and let
\(\E:X\to(-\infty,+\infty]\) be a proper convex functional with a unique
minimizer \(u^*\) in the admissible space. Its effective domain is
\[
  \operatorname{dom}\E=\{v\in X:\E(v)<+\infty\}.
\]

The descent arguments are carried out on an open neighborhood of the
relevant iterates where \(\E\) is finite and continuously Fr\'echet
differentiable. The energy space and topology of this neighborhood
are specified in the corresponding result. For energies with a restricted
effective domain, any additional topology is stated explicitly.
Arguments involving a Hessian further require twice Fr\'echet
differentiability on the stated neighborhood.

We denote the duality pairing between \(X^*\) and \(X\) by \(\ip rv\)
and define the energy defect by
\[
  \delta(u)=\E(u)-\E(u^*),\qquad \delta_m=\delta(G_m).
\]

\subsection{Ridge dictionaries and admissibility}

For an integer activation order \(k\geq1\), a raw ridge atom has the form
\[
  \varphi_{w,b}(x)=\relu(w\cdot x+b)^k,
  \qquad (w,b)\in\R^d\times\R.
\]
An admissible dictionary consists of nonzero ridge functions satisfying the constraints of \(X\),
after the corresponding linear projection. We use a symmetric normalized dictionary
\(\D_X\subset X\):
\[
g\in\D_X\Longrightarrow-g\in\D_X,\qquad \norm g_X=1\quad(g\in\D_X).
\]
When density is needed, we also require \(\overline{\spn\D_X}^{\,X}=X\).
We abbreviate \(\D_X\) to \(\D\) unless a different normalization is specified.

For pure-gradient Neumann energies the admissible space is
\[
  X_{p,\circ}=\left\{v\in W^{1,p}(\Omega):\int_\Omega v\dd x=0\right\},
\]
where \(\Omega\subset\mathbb R^d\) is the bounded, connected Lipschitz
domain considered in \Cref{sec:pgeometry}.
Each ridge is projected onto \(X_{p,\circ}\) before normalization. On a subdictionary with projected norms bounded below by a positive constant, uniform bounds for the projected atoms transfer to the normalized dictionary. The analytical dictionary consists of the nonzero projected atoms.

\begin{definition}[Atomic hull and atomic norm]
\label{def:atomic}
For a symmetric normalized dictionary \(\D\subset X\), define
\[
 K_1(\D)=\overline{\left\{\sum_{j=1}^Jc_jg_j:
 J<\infty,\ g_j\in\D,\ \sum_{j=1}^J\abs{c_j}\leq1\right\}}^{\,X},
\]
and
\[
 \norm{v}_{\mathcal A(\D)}=
 \inf\{B>0:v/B\in K_1(\D)\}.
\]
\end{definition}

For a bounded set \(K\subset X\), its \(n\)-th dyadic entropy number is
\[
 \varepsilon_n(K;X)=\inf\left\{h>0:\exists x_1,\ldots,x_{2^{n-1}}\in X,
 K\subset\bigcup_{j=1}^{2^{n-1}}(x_j+hB_X)\right\},\qquad n\geq1,
\]
where $B_X=\{x\in X:\norm x_X\leq1\}$. Thus the covering balls may have centers in $X$.
We write
$A\lesssim B$ when $A\leq CB$ with a constant independent of the iteration index; a
subscript on $\lesssim$ lists the permitted parameter dependence, and $A\simeq B$ denotes the two
corresponding inequalities.

\begin{assumption}[Atomic source]
\label[assumption]{ass:atomic}
There is a finite \(B_{\rm src}>0\), with \(B_{\rm src}\geq\norm{u^*}_{\mathcal A(\D)}\), such that \(u^*/B_{\rm src}\in K_1(\D)\).
\end{assumption}

The source condition uses the continuous dictionary normalized in \(X\). In one space dimension, \Cref{prop:1dsource} gives a sufficient condition for finite atomic norm.

\begin{proposition}[A one-dimensional \(\relu^k\) source class]
\label{prop:1dsource}
Let \(I=(0,1)\), \(k\geq1\), \(1\leq p<\infty\), and \(X=W^{1,p}(I)\). Suppose that the continuous
dictionary contains the normalized nonzero truncated powers \((x-t)_+^k/\|(\cdot-t)_+^k\|_X\),
\(0\leq t<1\), together with normalized fully active ridge functions spanning
\(\mathbb P_k\), the space of polynomials of degree at most \(k\) on \(I\). If
\(v\in W^{k+1,1}(I)\), then \(v\) belongs to the atomic space and
\[
\norm v_{\mathcal A(\D)}\leq C_{k,p}\left(
 \sum_{j=0}^k\abs{v^{(j)}(0)}+\norm{v^{(k+1)}}_{L^1(I)}\right),
\]
where \(C_{k,p}\) depends on the exponent \(p\), the interval, and the normalization of the raw atoms.
\end{proposition}

The proof in Supplement~S4 uses the integral remainder formula. Multivariate source and entropy estimates follow from shallow-network variation-space and width results under their domain and norm hypotheses \citep{SiegelXuVariation2023,SiegelXuSharp2024,MaoSiegelXu2026}.

\begin{lemma}[Projection and renormalization]
\label{lem:projection}
Let \(P:X\to X_0\) be bounded and linear. If a subdictionary \(\D_0\subset\D\) satisfies
\[
 0<c_0\leq\norm{Pg}_{X_0}\leq C_0<\infty\qquad(g\in\D_0),
\]
and \(\widetilde\D_0=\{Pg/\norm{Pg}_{X_0}:g\in\D_0\}\), then
\[
 P K_1(\D_0)\subset C_0K_1(\widetilde\D_0).
\]
The same computation also shows that a bounded linear map increases every entropy radius by at most
\(\norm P\), while the inverse normalization loses at most the factor \(c_0^{-1}\):
\[
 K_1(\widetilde\D_0)\subset c_0^{-1}P K_1(\D_0).
\]
\end{lemma}

\begin{proof}[Proof sketch]
For each atomic combination, write $Pg=\|Pg\|_{X_0}(Pg/\|Pg\|_{X_0})$ and use the upper and lower bounds on $\|Pg\|_{X_0}$.  Mapping a cover through $P$ multiplies its radius by at most $\|P\|$, while renormalization costs at most $c_0^{-1}$.
\end{proof}

\begin{corollary}[Transfer of source and entropy inputs]
\label{cor:projected-source-entropy}
In the setting of \Cref{lem:projection}, suppose \(u^*=Pu^\dagger\),
\(u^\dagger/B\in K_1(\D_0)\), and
\[
 \varepsilon_n(K_1(\D_0);X)\leq C_{\rm ent}n^{-\beta}.
\]
Then
\[
 \frac{u^*}{BC_0}\in K_1(\widetilde\D_0),
 \qquad
 \varepsilon_n(K_1(\widetilde\D_0);X_0)
 \leq c_0^{-1}\norm P\,C_{\rm ent}n^{-\beta}.
\]
Thus a ReLU source/convex-hull entropy theorem on an extended parameter domain passes to the projected,
renormalized feasible dictionary whenever the nondegeneracy bounds in \Cref{lem:projection} hold.
\end{corollary}

\begin{proof}[Proof sketch]
Apply the two inclusions in \Cref{lem:projection} to the atomic representation of $u^\dagger$ and to an entropy cover of $K_1(\D_0)$.  The source radius gains the factor $C_0$, and the covering radius gains $c_0^{-1}\|P\|$.
\end{proof}

\subsection{Dictionary residuals}

For \(r\in X^*\), introduce the dictionary dual seminorm
\[
  \norm r_{\D^*}=\sup_{g\in\D_X}\abs{\ip r g}.
\]
This quantity measures the derivative in the available dictionary directions. It is a seminorm;
if the dictionary has dense linear span, its vanishing implies that the derivative is zero.

\subsection{Ideal and finite-pool Chebyshev greedy algorithms}

When the derivative-score supremum is attained, the continuous CGA
selects a maximizing atom and minimizes the energy over the enlarged
trial space. Requiring only a prescribed fraction \(t_m\in(0,1]\) of
the supremum gives the weak CGA (wCGA), described in
\Cref{alg:continuous}.

\begin{algorithm}[htbp]
\caption{Continuous weak Chebyshev greedy algorithm}
\label{alg:continuous}
Assume \(\E(0)<\infty\), choose weakness numbers \(t_m\in(0,1]\), and set
\(G_0=0\), \(\Phi_0=\{0\}\). Given
\(G_m\) and \(\Phi_m=\spn\{g_1,\ldots,g_m\}\):
\begin{enumerate}[label=\textup{(\roman*)}]
  \item if \(\norm{\E'(G_m)}_{\D^*}=0\), stop; otherwise choose \(g_{m+1}\in\D_X\) with descent sign such that
  \[
  -\ip{\E'(G_m)}{g_{m+1}}
  \geq t_m\norm{\E'(G_m)}_{\D^*};
  \]
  \item set \(\Phi_{m+1}=\spn(\Phi_m,g_{m+1})\) and compute
  \[
    G_{m+1}=\argmin_{v\in\Phi_{m+1}}\E(v).
  \]
\end{enumerate}
\end{algorithm}

The choice \(t_m=1\) corresponds to exact selection, available, for
example, on a compact nondegenerate parameter set with continuous
score. Since changing the sign of an atom leaves its linear span
unchanged, the selected atom can be assigned its descent sign.

Exact reoptimization yields the variational orthogonality
\begin{equation}
 \ip{\E'(G_m)}v=0\qquad(v\in\Phi_m).
 \label{eq:orthogonality}
\end{equation}
The finite-pool algorithm uses a finite admissible pool \(\Ppool_m\)
and, at the current computed state \(\widehat G_m\), a numerically
evaluated score \(\widetilde s_m:\Ppool_m\to[0,\infty)\) approximating
\(g\mapsto|\ip{\E'(\widehat G_m)}g|\). Coefficient minimization is
performed to a finite tolerance.

\begin{algorithm}[htbp]
\caption{Finite-pool inexact CGA}
\label{alg:finite}
Set \(\widehat G_0=0\), \(\widehat\Phi_0=\{0\}\). At step \(m\), form a finite admissible pool
\(\Ppool_m\subset\D\). Let
\[
 s_m(g)=\abs{\ip{\E'(\widehat G_m)}g},
 \qquad \widetilde s_m(g)\geq0.
\]
Choose \(t_m^{\rm pool}\in(0,1]\) and \(\widehat g_{m+1}\in\Ppool_m\) satisfying
\[
 \widetilde s_m(\widehat g_{m+1})
 \geq t_m^{\rm pool}\max_{g\in\Ppool_m}\widetilde s_m(g),
\]
give it the true descent sign, and define
\(\widehat\Phi_{m+1}=\spn(\widehat\Phi_m,\widehat g_{m+1})\). Compute a coefficient vector whose state
obeys
\[
 \E(\widehat G_{m+1})
 \leq\inf_{v\in\widehat\Phi_{m+1}}\E(v)+\varepsilon_m^{\rm opt},
 \qquad \varepsilon_m^{\rm opt}\geq0.
\]
\end{algorithm}

For each run, the implementation records the reason for termination.

\subsection{Finite-pool error quantities}
\label{sec:controlled-selection}

The score error and continuous-dictionary coverage are quantified by
\begin{align}
 \sup_{g\in\Ppool_m}\abs{\widetilde s_m(g)-s_m(g)}&\leq\xi_m,
 \label{eq:score-error}\\
 \tau_m&=\frac{\max_{g\in\Ppool_m}s_m(g)}{\sup_{g\in\D}s_m(g)},\qquad 0\leq\tau_m\leq1,
 \label{eq:pool-coverage}
\end{align}
with \(\tau_m=1\) when the denominator is zero. The convergence estimates specify a positive lower bound on the effective selection strength \(t_m^{\rm pool}\tau_m\).

\begin{lemma}[True score selected from an evaluated pool]
\label{lem:true-score}
Under \cref{eq:score-error,eq:pool-coverage},
\begin{equation}
 s_m(\widehat g_{m+1})\geq
 t_m^{\rm pool}\tau_m\norm{\E'(\widehat G_m)}_{\D^*}
 -(1+t_m^{\rm pool})\xi_m.
\label{eq:true-score-lower}
\end{equation}
\end{lemma}

\begin{proof}[Proof sketch]
Apply the score-error bound before and after the approximate maximization, and express the finite-pool maximum through \(\tau_m\) and the continuous score supremum.
\end{proof}

For coefficient optimization with the discrete energy \(\E_h\),
the following estimate relates the correction tolerance to the
discrepancy between \(\E_h\) and \(\E\) on a suitable sublevel set of
the current trial space \(\widehat\Phi_{m+1}\).

Let \(G_{m+1}^{\sharp}\) denote an exact minimizer of \(\E\) over
\(\widehat\Phi_{m+1}\). For an energy level \(\Lambda_m^E\) chosen so that both
\(G_{m+1}^{\sharp}\) and the computed state \(\widehat G_{m+1}\) are included, define the relevant sublevel
set and its active-space restriction by
\begin{equation}
 \mathcal L_m=\{v\in X:\E(v)\leq\Lambda_m^E\},
 \qquad \mathcal L_m^{\rm act}=\widehat\Phi_{m+1}\cap\mathcal L_m.
 \label{eq:active-sublevel}
\end{equation}
Thus \(\mathcal L_m^{\rm act}\) is the part of the current finite-dimensional
trial space on which the discrete energy \(\E_h\) and the variational
energy \(\E\) are compared; in particular,
\(\inf_{v\in\mathcal L_m^{\rm act}}\E(v)=\E(G_{m+1}^{\sharp})\).

\begin{lemma}[Transfer from the discrete energy to the analysis energy]
\label{lem:energy-transfer}
Suppose on the active-space sublevel set \(\mathcal L_m^{\rm act}\) in
\cref{eq:active-sublevel} that
\[
 \sup_{v\in\mathcal L_m^{\rm act}}
 \abs{\E_h(v)-\E(v)}\leq\zeta_m^E.
\]
If the computed state satisfies
\[
 \E_h(\widehat G_{m+1})\leq
 \inf_{v\in\mathcal L_m^{\rm act}}\E_h(v)+\varepsilon_m^{\rm opt,h},
\]
then \Cref{alg:finite} holds with
\(\varepsilon_m^{\rm opt}\leq\varepsilon_m^{\rm opt,h}+2\zeta_m^E\).
\end{lemma}

\begin{proof}[Proof sketch]
Compare the computed state with the exact active minimizer in the common sublevel set.  The uniform energy discrepancy costs at most $\zeta_m^E$ at each endpoint, so the two discrepancies add to $2\zeta_m^E$.
\end{proof}
